\documentclass[10pt,a4paper]{amsart}
\usepackage[margin=19mm]{geometry}
\usepackage{amsmath,amssymb,mathtools}
\usepackage[scr=rsfs,cal=euler]{mathalfa}
\usepackage{xfrac}
\usepackage[unicode,hidelinks,bookmarks=false]{hyperref}
\newcommand{\ie}{i.e.\ }
\newcommand{\cf}{cf.\ }
\newcommand{\resp}{respectively\ }
\newcommand{\Subsection}{\subsection}
\providecommand{\nobreakdash}{\nobreak\hbox{-}}
\setlength{\emergencystretch}{2em}
\multlinegap0pt
\usepackage{bm}
\usepackage{bbm}
\usepackage{dsfont}
\usepackage{xspace}
\usepackage{cancel}
\usepackage{mathtools}
\usepackage{amsthm}
\makeatletter
\@ifundefined{lem}{\newtheorem{lem}{Lem}}{}
\@ifundefined{thm}{\newtheorem{thm}{Theorem}}{}
\@ifundefined{cor}{\newtheorem{cor}[thm]{Corollary}}{}
\makeatother
\makeatletter
\newcommand{\A}{\mathcal A}
\DeclareMathOperator{\Ad}{Ad}
\newcommand{\B}{\mathcal B}
\newcommand{\BB}{\mathbb B}
\DeclareMathOperator{\Bis}{Bis}
\newcommand{\EE}{\mathcal E}
\newcommand{\G}{\mathcal G} %groupoid
\renewcommand{\L}{\mathcal L}
\newcommand{\M}{\mathcal M}
\newcommand{\Mat}{\mathbb M} % matrices
\newcommand{\N}{\mathbb N}
\newcommand*{\contb}{\mathrm{C_b}}%continuous functions vanishing at infinity
\newcommand*{\contc}{\mathrm{C_c}}%continuous functions with compact support
\newcommand*{\contz}{\mathrm{C_0}}%continuous functions vanishing at infinity
\newcommand{\cst}{\ifmmode\mathrm{C}^*\else{$\mathrm{C}^*$}\fi}
\DeclareMathOperator{\f}{f}% reduced
\newcommand{\m}{m}
\DeclareMathOperator{\rd}{r}% reduced
\DeclareMathOperator{\red}{red}% reduced
\expandafter\ifx\csname rlist\endcsname\relax
\newenvironment{rlist}
\fi
\makeatother
\title[Fourier--Stieltjes category for twisted groupoid actions]{Fourier--Stieltjes category for twisted groupoid actions}
\author{Alcides Buss, Bartosz Kwa\'sniewski, Andrew McKee, Adam Skalski}
\date{Source: arXiv:2405.15653v2 (CC BY 4.0)}
\begin{document}
\maketitle

\noindent\emph{Source and licence.} Fourier--Stieltjes category for twisted groupoid actions. Version arXiv:2405.15653v2 (CC BY 4.0), distributed under the Creative Commons Attribution 4.0 International licence, as recorded in the arXiv licence metadata for this exact version and shown on the abstract page https://arxiv.org/abs/2405.15653v2. Full rights notice: licenses/arxiv-2405.15653-CC-BY-4.0.txt.\par\smallskip

\noindent\textit{Editorial scope.} Counted unit: the Thm whose source label is \texttt{thm:FSmultiplier\_positive\_definite} in the article (source0.tex lines 2749--2767, source label \texttt{thm:FSmultiplier\_positive\_definite}), delivered together with the article's own complete original proof of it (source0.tex lines 2768--2903, one \texttt{proof} environment of 136 lines). No mathematics is written, repaired or re-derived: statement and proof are the article's own text line for line. Required context retained uncounted: enu:Murphy\_Stinespring2 (lines 738--752); cor:Murphy\_Stinespring (lines 832--850); thm:FSmultiplier\_extends\_to\_reduced\_and\_full (lines 2534--2538); lem:completely\_positive\_implies\_positive\_definite (lines 2664--2668); lem:positive\_action (lines 2702--2716). Every definition, display and label of those lines is retained unchanged. Omitted from this package and not redistributed: 1--737, 753--831, 851--2533, 2539--2663, 2669--2701, 2717--2748, 2904--4061. Nothing inside a retained excerpt is omitted or altered: every retained body is delivered exactly as the article's lines are written, so no omission receipt of any kind accompanies this package. Citations: the retained excerpts carry no citation command at all, so no bibliography entry of the article is transcribed and no reference is redistributed. Reference adaptation: none was needed and none was made; every reference inside the delivered unit resolves inside it, so no unresolved reference or citation is produced. Printed slips kept literal: no printed word, symbol, hypothesis or number of the counted statement or of its proof was altered. Layout and class: the article's own class is \texttt{amsart} with packages including xy, amsfonts, amsthm, amssymb, latexsym, amsmath, hyperref, comment, a4wide, color, xcolor, inputenc, fontenc, lmodern; the wrapper is the lane's pinned \texttt{amsart} editorial wrapper at 10pt on A4, which restates the theorem-like environments the retained text needs and adds the article's own macro definitions that the retained text needs (\texttt{\textbackslash A}, \texttt{\textbackslash Ad}, \texttt{\textbackslash B}, \texttt{\textbackslash BB}, \texttt{\textbackslash Bis}, \texttt{\textbackslash EE}, \texttt{\textbackslash G}, \texttt{\textbackslash L}, \texttt{\textbackslash M}, \texttt{\textbackslash Mat}, \texttt{\textbackslash N}, \texttt{\textbackslash contb}, \texttt{\textbackslash contc}, \texttt{\textbackslash contz}, \texttt{\textbackslash cst}, \texttt{\textbackslash f}, \texttt{\textbackslash m}, \texttt{\textbackslash rd}, \texttt{\textbackslash red}), with extra wrapper packages (bm, bbm, dsfont, xspace, cancel, mathtools). The delivered numbering is the unit's own. Assets: no retained span contains a figure environment, a graphics-inclusion command, a PDF-page inclusion, a table or a TikZ picture; no image file is copied into this package, the package declares no asset dependency and no image file is redistributed with it. Licence: version arXiv:2405.15653v2 of this article is distributed under the Creative Commons Attribution 4.0 International licence, as recorded in the arXiv licence metadata for this exact version and shown on the abstract page https://arxiv.org/abs/2405.15653v2. A full rights notice is delivered at licenses/arxiv-2405.15653-CC-BY-4.0.txt. Characters: every retained excerpt is pure ASCII, so the delivered engine, XeLaTeX with its default Latin Modern text font, reports no missing character.
\begin{enumerate}
    \item\label{enu:Murphy_Stinespring2} there exists an $A$-$B$-$\cst$-correspondence $E$  and a map $\zeta:S \to \L(E_B,E)$ such that 
    \[ 
        k(s,t)(a) =  \zeta(s)^* \varphi_E(a) \zeta(t),\qquad   s, t \in S,\  a \in A,
    \]
    where $\varphi_E:A\to \L(E)$ is the $^*$-homomorphism defining the left action of $A$ on $E$.
    \item\label{enu:Murphy_Stinespring1} $k$ is a \emph{positive-definite kernel} in the sense that for all $n \in \N$, $s_1,\ldots, s_n \in S$, $a_1, \ldots ,a_n\in A$ and $\xi_1, \ldots, \xi_n \in E_B$, we have
    \[
        \sum_{i,j=1}^n \langle \xi_i, k(s_i,s_j) (a_i^*a_j)\xi_j\rangle_B \geq 0,
    \]
    and in addition $k$ is \emph{strict} in the sense that for each $s\in S$ the map $k(s,s):A\to \L(E_B)$ is strict, i.e.\ for an approximate unit $\{e_i\}_{i \in I}$ in $A$, the net $\{k(s,s)(e_i)\}_{i \in I}$ converges strictly in $\L(E_B)$.
%\item\label{enu:Murphy_Stinespring3}
% there exists a $\cst$-correspondence $E$ from $A$ to $B$ and a map $\zeta:S \to M(E)$ such that for all $s, t \in S$ and $a \in A$ we have
%\[ k(s,t)(a) = \langle  \cdot \zeta(s),a\cdot \zeta(t) \rangle.\]
\end{enumerate}

\begin{cor}\label{cor:Murphy_Stinespring}
Let $k:S \times S \to \BB(A,B)$ be a map where $A$ and $B$ are $\cst$-algebras and $S$ is a set.
The following are equivalent:
\begin{enumerate}
    \item\label{enu:cor_Murphy_Stinespring2} there exists a $\cst$-correspondence $E$ from $A$ to $B$ and a map $\zeta:S \to \M(E)$ such that 
    \[
        k(s,t)(a) = \langle   \zeta(s), a\cdot \zeta(t) \rangle_{\M(B)}, \qquad s, t \in S,\ a \in A;
    \]

    \item\label{enu:cor_Murphy_Stinespring1} $k$ is a \emph{positive-definite kernel} in the sense that for all $n \in \N$, $s_1,\ldots, s_n \in S$ and $a_1, \ldots, a_n \in A$ the matrix
    \[
        [k(s_i,s_j) ( a_i^*a_j)]_{i,j=1}^n \in \Mat_n(B)
    \]
    is positive, %equivalently, for all $n \in \N$, $s_1,\ldots, s_n \in S$ and $a_1, \ldots, a_n\in A, b_1,\ldots, b_n \in B$ we have \(\sum_{i,j=1}^n b_i^* k(s_i,s_j) (a_i^*a_j)b_j \geq 0.\)
    and $k$ is \emph{strict} in the sense that for each $s\in S$ the map $k(s,s):A\to B$ is strict, i.e.\ for an approximate unit $\{e_i\}_{i \in I}$ in $A$, the net $\{k(s,s)(e_i)\}_{i \in I}$ converges strictly in $\M(B)$.
\end{enumerate}
Moreover, if the conditions above hold, then the $\cst$-correspondence $E$ in \ref{enu:Murphy_Stinespring2} may be chosen so that the set $\{a\zeta(s)b:s \in S, a\in A, b\in B\}$ is dense in $E$; we call such pairs $(\zeta,E)$ \emph{minimal}. 
Given another minimal pair $(\zeta',E')$ as above, we have a uniquely determined unitary isomorphism of $\cst$-correspondences $U:E \to E'$ such that $U\zeta(s) = \zeta'(s)$, $s \in S$.
\end{cor}

\begin{thm}\label{thm:FSmultiplier_extends_to_reduced_and_full} 
Let $\A$, $\B$  be $\cst$-bundles over $X$, carrying twisted actions $(\alpha,u_\alpha)$ and $(\beta,u_\beta)$ of $\G$.
Every Fourier--Stieltjes multiplier $T\in FS[(\alpha,u_{\alpha}),(\beta,u_{\beta})]$ yields strict completely bounded maps $\m_T^{\rd} : \cst_{\red}(\A^{(\alpha,u_{\alpha})})\to \cst_{\red}(\B^{(\beta,u_{\beta})})$ and $\m_T^{\f} : \cst(\A^{(\alpha,u_{\alpha})})\to \cst(\B^{(\beta,u_{\beta})})$ between reduced and full crossed products, and $\|\m_T^{\rd}\|_{cb} , \|\m_T^{\f}\|_{cb} \leq \|T\|_{FS}$. 
If $T=T_{\EE,L,\xi,\xi}$ for some equivariant representation $L$ of $\G$ on a $\cst$-$\A$-$\B$-correspondence bundle $\EE$ and section $\xi \in \contb(\M(\EE))$, then $\m_T^{r}$ and $\m_T^{f}$ are completely positive and $\|\m_T^{\rd}\|_{cb} = \| \m_T^{\f} \|_{cb} = \| T \|_{FS} = \|\xi\|^2$.
\end{thm}

\begin{lem}\label{lem:completely_positive_implies_positive_definite}
Let $\A = \{ A_g \}_{g \in \G}$ and $\B = \{ B_g \}_{g \in \G}$ be Fell bundles over $\G$. 
Let $T = \{ T_g \}_{g\in \G}$ be a reduced or full multiplier from $\A$ to $\B$. 
If the corresponding map $\m_T$ between cross-sectional $\cst$-algebras is completely positive, then $T$ is positive-definite and bounded (in fact we have $\|\m_T\|=\sup_{x\in X}\|T_x\|$), and $T$ is strict if $\m_T$ is strict.
\end{lem}

\begin{lem}\label{lem:positive_action}
Let $\A$, $\B$ be $\cst$-bundles over $X$, carrying twisted actions $(\alpha,u_\alpha)$ and $(\beta,u_\beta)$ of $\G$.
For a  multiplier $T=\{T_g\}_{g\in \G}$ from $\A^{(\alpha,u_{\alpha})}$ to $\B^{(\beta,u_{\beta})}$ the following are equivalent:
\begin{enumerate}
    \item\label{enu:positive_action1} $T$ is positive-definite;
    \item\label{enu:positive_action2} for any $x \in X$, $g_1 , \ldots , g_n \in \G_x$, $n \in \N$, and any collection $\{ a_{g_i} \in A_x : 1 \leq i \leq n \}$ the following matrix is positive: 
	\[
	   \left( \beta_{g_i}^{-1} \left(T_{g_i g_j^{-1}} \big( \alpha_{g_i}(a_{g_i} a_{g_j}^*) u_{\alpha}(g_i g_j^{-1} , g_j)^* \big) u_{\beta}(g_i g_j^{-1} , g_j) \right) \right)_{i,j = 1}^n \in \Mat_{n}(B_x);
	\]
	\item\label{enu:positive_action3} for any $x \in X$, $g_1 , \ldots , g_n \in \G^x$, $n \in \N$, and any collection $\{ a_{g_i} \in A_x : 1 \leq i \leq n \}$ the following matrix is positive: 
	\[
	   \left( \beta_{g_i} \left(T_{g_i^{-1} g_j} \big( \alpha_{g_i}^{-1}(a_{g_i}^* a_{g_j} u_{\alpha}(g_i , g_i^{-1} g_j)^*) \big) \right) u_{\beta}(g_i , g_i^{-1} g_j) \right)_{i,j = 1}^n \in \Mat_{n}(B_x). 
	\]
	\end{enumerate}
\end{lem}

\begin{thm}\label{thm:FSmultiplier_positive_definite}
Let $\A$, $\B$ be $\cst$-bundles over $X$, carrying twisted actions $(\alpha,u_\alpha)$ and $(\beta,u_\beta)$ of $\G$, and set $A:=\contz(\A),$ $B:=\contz(\B)$. 
For any multiplier $T$ from $\A^{(\alpha,u_{\alpha})}$ to $\B^{(\beta,u_{\beta})}$, the following are equivalent: 
\begin{enumerate}
	\item\label{enu:FSmultiplier_positive_definite1} $T$ is  strict, bounded and positive-definite;
    \item\label{enu:FSmultiplier_positive_definite2} $\m_T$ %is a reduced, strict completely positive  multiplier, i.e. it 
	extends to a strict completely positive map $\m_T^{\rd} : \cst_{\red}(\A^{(\alpha,u_{\alpha})}) \to \cst_{\red}(\B^{(\beta,u_{\beta})})$;
	\item\label{enu:FSmultiplier_positive_definite3} $\m_T$ %is a full, strict completely positive multiplier, i.e. it 
	extends to a strict completely positive map $\m_T^{\f}:\cst(\A^{(\alpha,u_{\alpha})}) \to \cst(\B^{(\beta,u_{\beta})})$;
	\item\label{enu:FSmultiplier_positive_definite4} $T = T_{\EE,L,\xi,\xi}$ %\in FS[(\alpha,u_{\alpha}),(\beta,u_{\beta})]$  
	for some equivariant representation $L$ of $\G$ on a $\cst$-$\A$-$\B$-correspondence bundle $\EE$ and section $\xi\in \contb(\M(\EE))$.
\end{enumerate}	
If the above equivalent conditions hold, then  $\| \m_T^{\f} \|_{cb} = \| \m_T^{\rd} \|_{cb} = \| \m_T^{\f} \| = \| \m_T^{\rd} \| = \| T \|_{FS} = \| \xi \|^2$ and the triple $(\EE,L,\xi)$ in \ref{enu:FSmultiplier_positive_definite4} can be chosen to be cyclic, in the sense that the set
\[
    \{ L_U (a \xi b) : a \in \contz(\A|_{s(U)}) ,\ b \in \contz(\B|_{s(U)}) ,\ U\in \Bis(\G) \}
\]
is linearly dense in $\contz(\EE)$. 
Moreover, a cyclic triple $(\EE,L,\xi)$ in \ref{enu:FSmultiplier_positive_definite4} is unique up to a unitary equivalence in the sense that if  $(\EE',L',\xi')$ is another cyclic triple with $T=T_{\EE',L',\xi',\xi'}$, then there is an $A$-$B$-bimodule $\contz(X)$-unitary $W:\contz(\EE)\to \contz (\EE')$ establishing equivalence between $L$ and $L'$, and sending $\xi$ to $\xi'$.
\end{thm}

\begin{proof} 
The implications \ref{enu:FSmultiplier_positive_definite4}$\implies$\ref{enu:FSmultiplier_positive_definite2},\ref{enu:FSmultiplier_positive_definite3} follow from Theorem~\ref{thm:FSmultiplier_extends_to_reduced_and_full} and the implications \ref{enu:FSmultiplier_positive_definite2},\ref{enu:FSmultiplier_positive_definite3}$\implies$\ref{enu:FSmultiplier_positive_definite1} are a consequence of Lemma~\ref{lem:completely_positive_implies_positive_definite}.

Thus we only need to prove that \ref{enu:FSmultiplier_positive_definite1} implies \ref{enu:FSmultiplier_positive_definite4}. 
Fix a strict, bounded  and positive-definite multiplier $T$. 
For each $x \in X$ define $k_x : \G^x \times \G^x \to \B(A_x,B_x)$ via
\[
    k_x(g, h)(a) := \beta_{g} \left(T_{g^{-1} h} \big( \alpha_g^{-1} (au_{\alpha}(g,g^{-1}h)^*) \big) \right) u_{\beta} (g,g^{-1}h) % \beta_{g}^{-1} \left(T_{g h^{-1}} \big(\alpha_g(a) u_{\alpha}(gh^{-1},h)^*\big) u_{\beta}(gh^{-1},h)\right)
    \qquad g, h \in \G_x. 
\]
Then by Lemma~\ref{lem:positive_action}, $k_x$ is positive-definite and strict in the sense of Corollary~\ref{cor:Murphy_Stinespring}. 
Thus, by this latter result, there is a $\cst$-correspondence $E_x$ from $A_x$ to $B_x$ and a map $\xi_x : \G^x \to \M(E_x)$ such that the set $E_x^0 := \{ a \xi_x(g)b : g \in \G^x ,\ a \in A_x, b \in B_x \}$ is linearly dense in $E_x$, and for $g, h \in \G^x$, $a \in A_x$, we have 
\begin{equation} \label{kernelformula}
    \langle \xi_x(g), a \cdot \xi_x(h) \rangle_{\M(B_x)} =  \beta_{g} \left(T_{g^{-1} h} \big(\alpha_g^{-1}(au_{\alpha}(g,g^{-1}h)^*) \big) \right)u_{\beta}(g,g^{-1}h) \in B_x. 
\end{equation}
In this way, we get a family of $\cst$-correspondences $\EE:=\{E_x\}_{x\in X}$. 
Define the section $\xi : \G \to \M(\EE) = \{ \M(E_x) \}_{x\in X}$ by setting $\xi(g):=\xi_{r(g)}(g)$, $g \in \G$.
To topologise $\EE$ we let $\Gamma$ be the linear span of sections of the form $a\xi_U b$, where $U\in \Bis(\G)$, $a\in \contc(\A|_{r(U)})$, $b\in \contc(\B|_{r(U)})$ and
\[
    (a\xi_U b)(x) := a(x) \xi(r|_{U}^{-1}(x))b(x) , \qquad x\in X.
\]
If $c\xi_V d\in \Gamma$ is another element of this form, with $V\in \Bis(\G)$, $c\in \contc(\A|_{r(V})$, $d\in \contc (\B|_{r(V)})$, then $\langle a\xi_U b, c \xi_V d\rangle_{B} (x):=  \langle a\xi_U b(x), c \xi_V d (x)\rangle_{B_x}$ is zero outside $r(U)\cap r(V)$ and writing $g := r|_{U}^{-1}(x)$ and $h := r|_{V}^{-1}(x)$ for $x\in r(U)\cap r(V)$, by \eqref{kernelformula} we get 
\[
    \langle a\xi_U b, c \xi_V d\rangle_{B} (x) = b(x)^* \Big( \beta_{g} \Big(T_{g^{-1} h} \big( \alpha_g^{-1}( a(x)^* c(x)u_{\alpha}(g,g^{-1}h)^*) \big) \Big)u_{\beta}(g,g^{-1}h)\Big) d(x). 
\]
By continuity of all the operations involved, we see that the map $X\ni x\mapsto \langle a\xi_U b, c \xi_V d\rangle_{B} (x) \in B_x$ belongs to $\contc(\B)$. 
It extends uniquely to the (positive-definite) conjugate-linear map $\langle \cdot, \cdot \rangle_{B}:\Gamma\times \Gamma \to \contc(\B)\subseteq B$. 
In particular, for any $a\in \Gamma$ the map $X\mapsto \|a(x)\|^2=\|\langle a, a \rangle_{B}(x)\|$ is upper semicontinuous.
The fact that for each $x \in X$ the values of the prescribed sections are dense in the fibre $E_x$ follows immediately from the minimality assumption on each $E_x$ and the above construction of $\Gamma$.
Hence $\EE$ admits a unique topology making it a Banach bundle  such that $\Gamma\subseteq \contc(\EE)$.
The left $\A$ and right $\B$-module actions on $\EE$ are continuous by the very definition of sections, and the $B$-valued inner product $\langle \cdot, \cdot \rangle_{B}$ extends uniquely to $\contz(\EE)$ so that $\contz(\EE)$ becomes a $\contz(X)$-$\cst$-correspondence from $A$ to $B$. 

We claim that there is an $(\alpha,u_{\alpha})$-$(\beta,u_{\beta})$-equivariant representation of $\G$ on $\EE$ determined by
\begin{equation}\label{eq:Ldef} 
    L_g (a \xi(h)b) := \alpha_g(a)u_{\alpha}{(g,h)} \xi(gh)u_{\beta}(g,h)^*\beta_g(b),
\end{equation}
where $g\in \G$, $h \in \G^{s(g)}, a \in A_{s(g)}, b \in B_{s(g)}$. 
The most difficult (technically involved) part of the proof is showing that for $h, f \in \G^{s(g)}$, $a, c \in A_{s(g)}$, $b, d \in B_{s(g)}$, we have 
\begin{equation}\label{eq:twisted_thing_to_prove} 
    \langle L_g (a \xi(h)b) , L_g (c \xi(f)d) \rangle_{B_{r(g)}} = \beta_g(\langle a \xi(h)b , c \xi(f)d\rangle_{B_{s(g)}}).
\end{equation}
We denote by $L$ and $R$ the left and right hand side of \eqref{eq:twisted_thing_to_prove}. 
To make the calculations easier to read and write we abbreviate $u_{g,h} := u_{\alpha}{(g,h)}$ and $u_{g,h}:=u_{\beta}{(g,h)}$ as it is clear from the context whether we twist $\alpha$ or $\beta$. 
Then 
\[
\begin{split}
    R &= \beta_g \Big( b^*\beta_{h} \Big(T_{h^{-1} f} \big( \alpha_h^{-1}( a^* c\cdot  u^{*}_{h,h^{-1}f} \big) \Big) u_{h,h^{-1}f}\, d\Big) \\
    &= \beta_g(b^*)u_{g,h}\,\beta_{gh} \Big( T_{h^{-1} f} \big( \alpha_h^{-1}( a^* c\cdot  u^{*}_{h,h^{-1}f}) \big) \Big) u^{*}_{g,h} \beta_g(u^{\beta}_{h,h^{-1}f}) \beta_g(d) \\
    &= \beta_g(b^*)u_{g,h} \beta_{gh} \Big( T_{h^{-1} f} \big( \alpha_h^{-1}( a^* c\cdot  u^{*}_{h,h^{-1}f}) \big) \Big) u_{gh,h^{-1}f} u^{*}_{g,h} \beta_g(d) ,
\end{split}
\]
where in the last equality we used the cocycle identity. 
The left hand side is equal to 
\[
\begin{split} 
    L &= \Big \langle \alpha_g(a)u_{g,h} \xi(gh)u^{*}_{g,h}\beta_g(b) , \alpha_g(c)u^{}_{g,f} \xi(gf)u^{*}_{g,f}\beta_g(d)  \Big\rangle_{B_{r(g)}} \\
        &= \beta_g(b^*)u_{g,h}\beta_{gh} \Big(T_{h^{-1} f} \big(\alpha_{gh}^{-1}(u^{*}_{g,h}  \alpha_g(a^* c)u_{g,f} u^{*}_{gh,h^{-1}f}) \big)\Big) u_{gh,h^{-1}f}u^{*}_{g,h} \beta_g(d).
\end{split}
\]
Thus we need to check that the terms that $T_{h^{-1} f}$ acts on are equal. 
We compute 
\[
\begin{split} 
    \alpha_{gh}^{-1} \big(u^{*}_{g,h} \alpha_g(a^* c)u_{g,f} u^{*}_{gh,h^{-1}f}\big) &= \alpha_{h}^{-1} \alpha_{g}^{-1} \Ad_{u_{g,h}} \big( u^{*}_{g,h}  \alpha_g(a^* c)u_{g,f} u^{*}_{gh,h^{-1}f} \big) \\
        &= \alpha_{h}^{-1} \Big( \alpha_{g}^{-1} \big(  \alpha_g(a^* c)u_{g,f} u^{*}_{gh,h^{-1}f} u^{*}_{g,h} \big) \big) \\
        &= \alpha_{h}^{-1} \Big( a^* c\cdot \alpha_{g}^{-1} \big( u_{g,f} u^{*}_{gh,h^{-1}f} u^{*}_{g,h} \big) \Big) \\
        &= \alpha_{h}^{-1}( a^* c\cdot  u^{*}_{h,h^{-1}f} ) , 
\end{split}
\]
where the equality $\alpha_{g}^{-1}(u_{g,f} u^{*}_{gh,h^{-1}f} u^{*}_{g,h}) = u^{*}_{h,h^{-1}f}$ at the last step follows from the cocycle identity, as we have $ \alpha_g (u_{h,h^{-1}f})u_{g,f} = u_{g,h}u_{gh,h^{-1}f}$. 
This proves that $L=R$, that is \eqref{eq:twisted_thing_to_prove} holds.
This last equality implies that the linear extension of the prescription in \eqref{eq:Ldef} defines consistently a linear isometric map $L_g: E_{s(g)}^0 \to E_{r(g)}^0$ such that, for all $\xi_1, \xi_2 \in E_{s(g)}^0$, 
\begin{equation}\label{Lscalar}
    \langle L_h(\xi_1), L_h(\xi_2) \rangle_{B_{r(h)}} = \beta_h( \langle \xi_1, \xi_2 \rangle_{B_{s(h)}}) . 
\end{equation}
By continuity we get an isometry $L_g \in \B(E_{s(g)}, E_{r(g)})$ where \eqref{Lscalar} holds for all $\xi_1, \xi_2 \in E_{s(g)}$. 
Next note that \eqref{eq:Ldef} readily implies that 
\[ 
    L_g (a \cdot a_0 \xi(h)b_0\cdot b) = \alpha_g(a)\cdot  L_g(a_0 \xi(g)b_0) \cdot \beta_g(b), 
\]
for all $a,a_0 \in A_{s(g)}$, $b,b_0 \in B_{s(g)}$, $h\in \G^{s(g)}$. 
Thus by linearity and continuity 
\[ 
    L_g (a\cdot \xi \cdot b) = \alpha_g(a)\cdot L_g(\xi) \cdot \beta_g(b), \qquad  \xi \in E_{s(g)} ,\ a \in A_{s(g)} ,\ b \in  B_{s(g)}. 
\]
Similarly, to see that for $(f,g)\in \G^{(2)}$ we have $L_{f}\circ L_{g}(\cdot) = u_{\alpha}(f,g)L_{fg}(\cdot)u_{\beta}(f,g)^*$ on $E_{s(g)}$, we only need to check it on $a\xi(h)b$ where $f\in \G^{s(g)}$, $a \in A_{s(g)}$, $b \in B_{s(g)}$. 
We compute using our shorthand notation for twists: 
\[
\begin{split} 
    L_{f}( L_{g}(a\xi(h)b)) &= L_f\Big(\alpha_g(a)\,u_{g,h} \cdot\xi(gh)\cdot u_{g,h}^*\,\beta_g(b)\Big) \\
        &= \alpha_f\big(\alpha_g(a) u_{g,h}\big) u_{f,gh}\cdot \xi(fgh) \cdot u_{f,gh}^* \,\beta_f\big(u_{g,h}^* \beta_g(b)\big) \\
        &= \alpha_f\big(\alpha_g(a)\big) \big( \alpha_f (u_{g,h}) u_{f,gh}\big) \cdot \xi(fgh) \cdot \big( u_{f,gh}^* \,\beta_f(u_{g,h}^*)\big) \beta_f\big(\beta_g(b)\big) \\
        &= \alpha_f\big(\alpha_g(a)\big) (u_{f,g}  u_{fg,h}) \cdot \xi(fgh) \cdot (u_{fg,h}^* u_{f,g}^*) \beta_f\big(\beta_g(b)\big) \\
        &= u_{f,g}\alpha_{fg}(a)u_{fg,h}\cdot \xi(fgh) \cdot u_{fg,h}^*  \, \beta_{fg}(b) u_{f,g}^* \\
        &= u_{f,g} \, L_{fg}(a\xi(h)b) \, u_{f,g}^* . 
\end{split}
\]
Hence $L = \{ L_{g} \}_{g \in \G}$ is an $(\alpha,u_{\alpha})$-$(\beta,u_{\beta})$-equivariant representation of $\G$  on $\EE$.
Moreover, for any $g\in \G$ and $a \in A_{r(g)}$, \eqref{kernelformula} implies that $T_g(a) = \langle \xi(r(g)), a \xi(g) \rangle_{\M(B_{r(g)})}\in B_{r(g)}$ and \eqref{eq:Ldef} implies that $\overline{L}_{g}(\xi(s(g)) = \xi(g)$.
Thus
\[
    T_g(a) = \langle \xi(r(g)), a \overline{L}_g\big(\xi(s(g))\big)\rangle_{\M(B_{r(g)})}, 
\]
so that $T=T_{\EE,L,\xi|_{X},\xi|_{X}}$ where $\xi|_{X} \in \contb(\M(\EE))$ is bounded because $T$ is bounded, and in view of the above we have $\|\xi(x)\|^2=\|T_x\|$ for every $x\in X$. 
This finishes the proof of \ref{enu:FSmultiplier_positive_definite1}$\implies$\ref{enu:FSmultiplier_positive_definite4}. 
Moreover, in view of \eqref{eq:Ldef} we have $L_{U}(a\xi|_{X} b)=\alpha_{U}(a)\xi_{U}\beta_{U}(b)$ for every $a\in A_{s(U)}$, $b\in A_{s(U)}$, $U\in \Bis(\G)$. 
The linear span of such elements is equal to $\Gamma$, which by construction is dense in $\contz(\EE)$. 
Hence the triple $(\EE,L,\xi|_{X})$ is cyclic as in the second part of the assertion.

Finally, let $(\EE,L,\xi)$, $(\EE',L',\xi')$, be two cyclic triples such that $T_{\EE',L',\xi',\xi'}=T_{\EE,L,\xi,\xi}$.
We claim that for $a\in \contz(\A|_{s(U)})$,  $U\in \Bis(\G)$, and $c\in \contz(\A|_{s(V)})$,  $V\in \Bis(\G)$ 
\begin{equation}\label{eq:inner_product_preserve_cyclic}
    \langle  L_U(a \xi),  L_{V}(b\xi )\rangle = \langle  L_U'(a \xi' ),   L_{V}'(c\xi' )\rangle.
\end{equation}
To prove this we express the left hand side only in terms of $a,c$, $U,V$ and $T$.  
If $x\notin r(U)\cap r(V)$, then $\langle L_U (a\xi),L_V(b\xi) \rangle (x)=0$.
Assume then that $x\in r(U)\cap r(V)$. 
There are unique $g\in U$ and $h\in V$ such that $r(h)=r(g)=x$. 
Put $a_0=a(s(g))^* \alpha_g^{-1}(u_{\beta}(g, g^{-1}h)^*)$ and $c_0=\alpha_{g^{-1}{h}}(c(s(h)))$ and compute 
\[
\begin{split}
    \langle L_U (a\xi),L_V(c\xi) \rangle (x) &= \langle  L_{g}(a\xi(s(g))),L_{h}(c\xi(s(h)) \rangle \\
        &= \langle  L_{g}(a\xi(s(g))),L_{h}(c\xi(s(h))) u_{\beta}(g, g^{-1}h)^*\rangle u_{\beta}(g, g^{-1}h) \\
        &= \langle  L_{g}(a\xi(s(g))), u_{\beta}(g, g^{-1}h)^*L_{g} L_{g^{-1}h}(c\xi(s(h))) \rangle u_{\beta}(g, g^{-1}h) \\
        &= \beta_g\left(\langle  \xi(s(g)),  a(s(g))^* \alpha_g^{-1}(u_{\beta}(g, g^{-1}h)^*)L_{g^{-1}h}(c\xi(s(h))) \rangle\right) u_{\beta}(g, g^{-1}h) \\
        &= \beta_g\left( \langle \xi(s(g)), a_0 c_0L_{g^{-1}h}(\xi(s(h))) \rangle\right) u_{\beta}(g, g^{-1}h) \\
        &= \beta_g\left(T_{gh^{-1}}(a_0 c_0)\right) u_{\beta}(g, g^{-1}h).
\end{split}
\]
This proves \eqref{eq:inner_product_preserve_cyclic}, which implies that for $a\in \contz(\A|_{s(U)})$, $b\in \contz(\B|_{s(U)})$, $U\in \Bis(\G)$, and $c\in \contz(\A|_{s(V)})$, $d\in \contz(\B|_{s(V)})$, $V\in \Bis(\G)$ we have
\[
    \langle L_U(a \xi b), L_{V}(c\xi d)\rangle = \langle L_U'(a \xi' b) , L_{V}'(c\xi' d)\rangle.
\]
Thus if the triples $(\EE,L,\xi)$, $(\EE',L',\xi')$ are cyclic, then by linearity and continuity 
the formula $W (L_U(a \xi b)):=L_U'(a \xi' b)$ determines a unitary  $W:\contz(\EE)\to \contz (\EE')$ with the properties described in the last part of the assertion.
\end{proof} 

\end{document}
